% HD-F1-T10. GPT-6 Astra Pro, /root/transport; 2026-11-22 UTC.
\section{HD-F1-T10: elimination of the terminal tensor branch over
\texorpdfstring{\(\mathbf Q(\sqrt{-7})\)}{Q(sqrt(-7))}}
\label{sec:f1-terminal-tensor-seven}

In the exact terminal tensor alternative of HD-F1-T08--T09,
the field \(K=\mathbf Q(\sqrt{-7})\) does not occur.
Consequently that alternative remains only for
\(\mathbf Q(i)\) and \(\mathbf Q(\sqrt{-2})\), still with
\(r\ge20\), coefficient rank \(4r\ge80\), and \(1\le d_0\le4\).
The exclusion uses the actual quotient factor and does not identify
its six-dimensional orthogonal complement with the whole
rank-\(6+d_0\) coefficient factor.

\paragraph{1. A reflection in the actual rank-four coefficient group.}
Retain the actual finite group
\(G=\mathfrak S_s\times\Gamma\), its representation
\(\operatorname{Aug}(K^s)\otimes_K M\), and all quotient curves
and maps of (F1T9.2). The representation of \(\Gamma\) on \(M\)
is faithful and absolutely irreducible of rank four.
Its character field is exactly \(K\): the equality
\[
 \operatorname{tr}_{\operatorname{Aug}\otimes M}(1,\gamma)
       =(s-1)\operatorname{tr}_M(\gamma)
 \tag{F1T10.1}
\]
and the rational augmentation character give both inclusions
of the two character fields.

The actual \(\Gamma\)-cover \(X\to Y\) supplied by T09 has
\[
 D=\dim_K H^1_{\mathrm{par}}(Y,M)=6+d_0,\qquad
 7\le D\le10,\qquad
 D=4(2g(Y)-2)+\sum_{j=1}^t e'_j .
 \tag{F1T10.2}
\]
Here \(t\) counts its nontrivial coefficient inertias and
\(e'_j=\operatorname{rank}(\gamma'_j-1)\ge1\).
The six-dimensional target is the image of the actual projector
\[
 \Pi=e_u(1-s^{-1}f^*f_*),
 \qquad H^1_{\mathrm{par}}(Y,M)
     =f^*H^1_{\mathrm{par}}(\mathbf P^1,M)\perp\operatorname{im}\Pi.
 \tag{F1T10.3}
\]
The first summand has the base pairing multiplied by \(s\).
All these maps are flat maps of the actual local systems, so the
target period rank nine is bounded by the parameter dimension
of this very cover \(X\to Y\).

For \(g(Y)\ge3\), (F1T10.2) is impossible. For genus two,
\(\sum e'_j=D-8\le2\), and the pointed-base dimension is at most
\(3+t\le5\). Thus \(g(Y)=0\) or \(1\). At genus zero the
period bound gives \(t\ge12\), while
\(\sum e'_j=D+8\le18<2t\). At genus one it gives \(t\ge9\),
while \(\sum e'_j=D\le10<2t\); the unpointed case has
parameter dimension at most one and is already excluded.
In either case some \(e'_j=1\). Hence \(\Gamma\) contains an
actual reflection in its specified rank-four representation.

\paragraph{2. Exact reflection blocks over \(K=\mathbf Q(\sqrt{-7})\).}
Let \(W\) be the normal subgroup of \(\Gamma\) generated by all
its reflections. Average a positive \(K\)-Hermitian form over
\(\Gamma\), join nonorthogonal \(K\)-root lines, and let \(U\)
be the span of one connected component. As proved with the
literal maps in HD-G27, paragraph 1, these spans are mutually
orthogonal, absolutely irreducible reflection modules, transitively
permuted by \(\Gamma\). Their common dimension \(a\) divides four.
For their stabilizer \(H\), the actual coordinate map is
\[
 \bigoplus_{j\in\Gamma/H}U\longrightarrow M,\qquad
       (u_j)\longmapsto\sum_j\rho(t_j)u_j .
 \tag{F1T10.4}
\]
If \(\sigma:H\to GL_K(U)\) is the stabilizer representation,
its induced trace formula and its actual \(K\)-matrices give
\[
 K=\mathbf Q(\operatorname{tr}_M\Gamma)
     \subseteq\mathbf Q(\operatorname{tr}_U\sigma(H))
     \subseteq K.
 \tag{F1T10.5}
\]

No finite representation of rank at most two over this field
can have nonrational character. Indeed an eigenvalue \(\zeta_N\)
whose irreducible polynomial over \(K\) has nonrational
coefficients requires
\(K\subset\mathbf Q(\zeta_N)\); otherwise its polynomial is
the rational cyclotomic polynomial. The inclusion forces
\(7\mid N\), since \(7\) ramifies in \(K\) and cyclotomic
ramification is supported on primes dividing \(N\). The
irreducible polynomial then has degree
\[
 [K(\zeta_N):K]=\varphi(N)/2\ge3.
 \tag{F1T10.6}
\]
This is exactly the cyclotomic restriction used in HD-R12,
applied to the actual stabilizer matrices. In rank at most two
all characteristic polynomials, and hence all traces, are
rational. Equations (F1T10.5)--(F1T10.6) exclude \(a=1,2\);
therefore \(a=4\). Thus \(W\) itself is irreducible on the
original \(M\), and there is no smaller replacement module.

\paragraph{3. The actual rank-four root models and normalizers.}
Use the published reflection-group classification in precisely
the trace-restricted form proved in HD-G27, paragraphs 3--4.
For clarity, its rank-four restriction is recorded explicitly.
The imprimitive groups are \(G(\ell,p,4)\). The traces of
\[
 \operatorname{diag}(\zeta,\zeta^{-1},1,1),\qquad
 \operatorname{diag}(\zeta,\zeta,\zeta^{-2},1)
 \tag{F1T10.7}
\]
first force \(\ell\in\{1,2,3,4,6\}\), and then exclude
\(\ell=3,4,6\) over \(K\). The case \(\ell=1\) is reducible.
Thus only \(B_4\) and \(D_4\) remain from this list.
The nonreal primitive rank-four groups \(G_{29},G_{31}\)
contain \(iI_4\), with trace \(4i\); \(G_{32}\) has an
order-three reflection of trace \(3+\omega\). These traces
are outside \(K\). The real group \(H_4\) has the two-reflection
trace \(3/2+\sqrt5/2\), also outside \(K\). The possible
irreducible root groups are exactly
\[
                         A_4,\ B_4,\ D_4,\ F_4.
 \tag{F1T10.8}
\]

These are actual \(K\)-models, not only abstract group names.
Solve the finite intertwiner equations
\(T\rho(w)=\rho_{\mathrm{root}}(w)T\) over \(K\), as in
G27. A nonzero solution is invertible by absolute irreducibility;
the complex equivalence and scalar extension guarantee a nonzero
solution over \(K\). Its pullback Hermitian form is a positive
rational scalar times the given one. Thus this changes
coordinates on the actual coefficient module and on every
parabolic cochain by the specified \(K\)-linear map.

In a real root model, any \(g\in\Gamma\) normalizing \(W\)
has \(g^{-1}\bar g\) scalar, since \(g\) and \(\bar g\)
induce the same conjugation on the rational matrices of \(W\).
Consequently \(g=aR\), where \(R\) is a real orthogonal
root-line normalizer and \(|a|=1\). The real normalizer is
finite, because its root-line permutation has kernel
\(\{\pm I\}\); finite order of \(g\) then makes \(a\) a root
of unity. The actual real normalizer maps are those checked
in G27:

For \(A_4\), the five maximal four-cliques in the root-line
graph are the stars of the complete graph on five vertices.
Their permutation recovers the literal \(S_5\) action on
the augmentation module. For \(B_4\), the reflection
conjugacy classes of sizes four and twelve distinguish the
coordinate lines, giving signed permutations. For \(D_4\),
the ordered orthogonal long-root frames have successive
counts \(24,6,4,2\); \(W(F_4)\) attains the resulting bound
1152 and is the full real normalizer. All these maps are
rational.

For \(F_4\), put
\[
 \mathcal S=\{\pm e_j\}\cup
       \{\tfrac12(\pm e_1\pm e_2\pm e_3\pm e_4)\},
 \quad
 \mathcal L=\{\pm e_j\pm e_k:j<k\},
 \quad
 M_0=\begin{pmatrix}
 1&1&0&0\\1&-1&0&0\\0&0&1&1\\0&0&1&-1
 \end{pmatrix}.
 \tag{F1T10.9}
\]
The connected bipartite normalized-root graph distinguishes
the short and long parts up to exchange. Its part-preserving
normalizer is \(W(F_4)\), by the same frame count.
The identities \(M_0\mathcal S=\mathcal L\) and
\(\frac12M_0\mathcal L=\mathcal S\) give the exact other coset:
\[
 N_{O_4(\mathbf R)}(W(F_4))
        =W(F_4)\ \cup\ (M_0/\sqrt2)W(F_4).
 \tag{F1T10.10}
\]

For every rational representative \(R\), the condition
\(aR\in GL_4(K)\) forces \(a\in K\cap\mu_\infty=\{\pm1\}\).
All such matrices therefore have rational trace. In the other
\(F_4\) coset, write \(g=(a/\sqrt2)M_0w\). A nonzero
rational entry gives \(a/\sqrt2\in K\). Hence
\(a^2\in K\cap\mu_\infty=\{\pm1\}\).
The value \(a^2=1\) would put \(1/\sqrt2\) in \(K\); the
value \(a^2=-1\) would put \(\sqrt{-2}/2\) in \(K\).
Both are impossible for \(K=\mathbf Q(\sqrt{-7})\).
Thus the whole \(\Gamma\) has rational character, contrary
to (F1T10.1). This contradicts the reflection forced in
paragraph 1 and proves the exclusion.

\paragraph{Scope, sources, and the retained complement map.}
The classification input is Cohen, \emph{Finite complex
reflection groups}, Ann.\ Sci.\ ENS 9 (1976), pp.379--436:
Theorem (2.4), pp.386--387; (4.4)(i), p.399; table (4.16),
p.412; and Theorem (5.12), p.435. The actual trace and real
normalizer proofs used above are carried through their complete
proofs in HD-G27, paragraphs 1,3,4; no new classification is
claimed. The present consequence is the exclusion for the
rank-\(6+d_0\) quotient supplied by T09, where the previously
proved rank-six whole-factor conclusion alone would not
suffice.

The exact maps in (F1T10.3) remain the maps on the actual
Jacobians: \(f^*\), \(f_*\), the primitive corner \(e_u\),
and the integral multiple clearing the denominator of \(\Pi\).
Their Poincare-biextension identity is
\(\widehat{f_*}\lambda_Z=\lambda_X f^*\), with
\((f^*)^*\Theta_X=s\Theta_Z\); the six-dimensional image is
the same polarized factor up to the already proved isogeny
and rank-six norm multipliers. Their variational cycle and
Hodge maps are the induced flat pullback/Gysin maps.
The actual graph correspondences define their \(K_0\) and
GRR realizations, with the degree scalars already displayed;
Adams operations act on the resulting degree-\(2j\) class
by the factor \(k^j\). No assertion about the determinant
of the whole quotient replaces the determinant of its
base-orthogonal complement.
